\title{xLSTM: Extended Long Short-Term Memory}

\begin{abstract}
The foundational concepts of the Long Short-Term Memory (LSTM) architecture—constant error carousel and gating mechanisms—were introduced in the 1990s. Since their inception, LSTMs have demonstrated remarkable resilience and have driven significant advances in deep learning, notably as the original backbone of early Large Language Models (LLMs). Despite their success, the emergence of Transformer models, characterized by highly parallelizable self-attention, signaled a paradigm shift, enabling models to eclipse LSTMs in large-scale applications. This work explores a pivotal question: To what extent can language modeling be enhanced by scaling LSTM-based models to billions of parameters, employing contemporary techniques from modern LLMs while addressing LSTM shortcomings? We address this by first proposing exponential gating, alongside innovative normalization and stabilization strategies. Furthermore, we alter the classical LSTM memory design, resulting in the following variants: (i) sLSTM, which introduces a scalar memory, scalar updating mechanism, and a new method for memory integration; (ii) mLSTM, a fully parallelizable model built on matrix memory and a covariance-based update approach. By embedding these enhanced LSTM components within residual block frameworks, we create xLSTM blocks, which, when residually layered, form scalable xLSTM architectures. The adoption of exponential gating and restructured memory forms equips xLSTM architectures to rival, and in certain cases surpass, the performance and scalability of current leading Transformer and State Space Model implementations.\\
Code available at: \url{https://github.com/NX-AI/xlstm}
\end{abstract}

\section{Introduction}
\label{sec:intro}

The Long Short-Term Memory (LSTM) ideas~\citep{Hochreiter:91,Hochreiter:97nips,Hochreiter:97}, i.e., the constant error carousel and gating, were introduced to overcome the vanishing gradient problem of recurrent neural networks~\citep{Hochreiter:91,Hochreiter:00book}:

\begin{align}
\label{eq:lstm_idea}
  c_t \ = \ f_t \ c_{t-1} \ + \ 
          i_t \  z_t \ , \quad  
          h_t \ = \ o_t \ \psi({c_t}) \ . 
\end{align}

The constant error carousel involves incrementally adjusting the cell state $c_{t-1}$ (indicated in green) through the addition of cell inputs $z_t$, with the process being regulated by sigmoid gates (depicted in blue). The input gate $i_t$ along with the forget gate $f_t$ dictate how this update unfolds, whereas the output gate $o_t$ determines what is released from the memory cell, namely the hidden state $h_t$. Before the cell state is utilized, it undergoes normalization or squashing via $\psi$, after which the output gating produces the hidden state.

Long Short-Term Memory networks (LSTMs) have seen widespread adoption across a broad spectrum of fields \citep{Hochreiter:01,Hochreiter:07,Schmidhuber:15}, and established dominance in text generation prior to the advent of Transformer architectures in 2017~\citep{Vaswani:17}. Their capability has been empirically validated in a variety of sequential problems, including but not limited to text production~\citep{Graves:13,Karpathy:15blog}, handwriting synthesis~\citep{Graves:13}, machine translation in sequence-to-sequence frameworks~\citep{Sutskever:14nips}, automated program evaluation~\citep{Zaremba:14arxiv}, image captioning~\citep{Karpathy:15,Hossain:19}, code generation~\citep{Karpathy:15blog}, modeling rainfall-runoff processes~\citep{Kratzert:18,Kratzert:19}, and hydrological prediction for flood warning systems~\citep{Nearing:24}. Within the reinforcement learning arena, LSTMs have set the standard for sequential modeling, serving as the backbone for prominent systems such as AlphaStar for StarCraft II~\citep{Vinyals:17short}, OpenAI Five for Dota 2~\citep{OpenAI:19}, and controllers for nuclear fusion magnetic devices~\citep{Degrave:22short}. LSTMs possess notable strengths in learning high-level abstractions, proficiently capturing semantic content in their memory mechanisms~\citep{Karpathy:15blog}, as demonstrated through the identification of neurons specializing in number and syntax~\citep{Lakretz:19}, linguistic features~\citep{Bau:19}, and sentiment~\citep{Radford:17arxiv}. Despite rapid advances in sequence modeling, LSTMs continue to play a vital role in contemporary applications~\citep{Degrave:22short,Nearing:24} and have proven their durability over time.

Although LSTMs have achieved remarkable results, they exhibit three principal shortcomings: (i) They cannot revise previously stored information. This issue is illustrated with the {\it Nearest Neighbor Search} problem (for details, see Appendix~\ref{sec:appNearestNeighborSearch}): Given a reference vector, the task involves sequentially scanning a sequence to locate the most similar vector and output the associated value at the end of the sequence. As depicted in the left panel of Figure~\ref{fig:lstmProblems}, mean squared error is reported for this task. LSTMs have difficulty replacing an already stored value when encountering a closer match later in the sequence, whereas our xLSTM overcomes this constraint by employing exponential gating mechanisms. (ii) They possess restricted memory capacity, requiring all information to be compressed into scalar cell states. This is demonstrated using {\it Rare Token Prediction}, as shown in the right panel of Figure~\ref{fig:lstmProblems}, where token prediction perplexity is plotted for subsets of tokens with varying frequencies on Wikitext-103~\citep{Merity:17}. LSTMs particularly underperform for rare tokens due to these storage limits, while xLSTM addresses this limitation by incorporating a matrix-based memory. (iii) Parallelization is hindered in LSTMs because of memory mixing, that is, the interconnections between hidden states at successive time steps, necessitating sequential computations.

These drawbacks have contributed to the rise of Transformer architectures~\citep{Vaswani:17} within the field of language modeling. A natural question arises: What levels of performance in language modeling can be realized if these limitations are surmounted and LSTMs are scaled to match the parameter counts of contemporary Large Language Models?

\section{Extended Long Short-Term Memory}
\label{sec:xLSTM}

To address the constraints present in standard LSTMs, Extended Long Short-Term Memory (xLSTM) incorporates two significant modifications as alternatives to the concept underlying Equation~\eqref{eq:lstm_idea}. These enhancements—namely, exponential gating mechanisms and the introduction of alternative memory structures—augment the LSTM framework, leading to the emergence of two new variants: (i) sLSTM (elaborated in Section~\ref{sec:sLSTM}), which is characterized by a scalar memory, scalar update, and incorporates mechanisms for memory mixing; and (ii) mLSTM (see Section~\ref{sec:mLSTM}), which utilizes a matrix-based memory system combined with a covariance (outer product) update, allowing for complete parallelization. Exponential gating is foundational to the improvements seen in both sLSTM and mLSTM models. For efficient parallel computation, mLSTM discards memory mixing (eliminating hidden-to-hidden recurrent links). Additionally, both variants are adaptable to the use of multiple memory cells: sLSTM allows for cell-wise memory mixing, while mLSTM treats multi-head and multi-cell configurations as functionally equivalent. sLSTM can also be extended to accommodate multiple heads, in which case, memory mixing occurs only among cells within each individual head and not across different heads. The combination of multi-head architecture and exponential gating in sLSTM thus introduces a novel form of memory interaction.

Incorporating these LSTM variants as foundational components within residual modules produces xLSTM blocks (refer to Section~\ref{sec:xLSTMblock}). Building deep models by stacking these xLSTM blocks residually gives rise to the overall xLSTM architectures (details in Section~\ref{sec:xLSTMarchitecure}). The layout of the xLSTM and its main constituents is depicted in Figure~\ref{fig:xlstm_sketch}.

\subsection{Review of the Long Short-Term Memory}
\label{sec:orgLSTM}

The foundational concept of the LSTM~\citep{Hochreiter:91,Hochreiter:97nips,Hochreiter:97} centers on a scalar memory cell designed for effective computation and information retention, thereby overcoming the vanishing gradient problem~\citep{Hochreiter:91,Hochreiter:00book} by means of the constant error carousel, which regulates the cell state transitions. The architecture includes three distinct gates: the input gate, the output gate, and the forget gate—the latter being subsequently incorporated by \citet{Gers:00}. The following update equations govern the behavior of the LSTM memory cell at time $t$:

\begin{align}
c_t \ &= \  f_t \ c_{t-1} \ + \ i_t \ z_t &  & &\text{cell state} \\
h_t  \ &= \ o_t \ \tilde{h}_t \ , & \tilde{h}_t \ &= \ \psi \left( c_t\right)
  &\text{hidden state} \\
z_t \ &= \ \varphi \left( \tilde{z}_t \right) \ , 
  &\tilde{z}_t \ &=  \ \mathbf{w}^\top_{z} \ \mathbf{x}_t \ + \
  r_{z}  h_{t-1} \ + \  b_{z} \ \
  &\text{cell input} \\
i_t \ &= \ \sigma \left( \tilde{i}_t  \right) \ , 
  &\tilde{i}_t \ &= \ \mathbf{w}^\top_{i} \ \mathbf{x}_t \ + \
  r_{i}  \ h_{t-1} \ + \  b_{i} \ \
  &\text{input gate} \\
f_t \ &= \ \sigma \left(  \tilde{f}_t \right) \ , 
  &\tilde{f}_t \ &= \ \mathbf{w}^\top_{f} \ \mathbf{x}_t  \ + \
  r_{f}  \ h_{t-1} \ + \  b_{f} \ \
  &\text{forget gate} \\
o_t \ &= \ \sigma \left( \tilde{o}_t \right) \ , 
  &\tilde{o}_t  \ &= \ \mathbf{w}^\top_{o} \ \mathbf{x}_t \ + \
  r_{o}  \ h_{t-1} \ + \  b_{o} \ \
  &\text{output gate} 
\end{align}

The weight vectors $\mathbf{w}_{z}$, $\mathbf{w}_{i}$, $\mathbf{w}_{f}$, and $\mathbf{w}_{o}$ serve as the input weight parameters linking the input vector $\mathbf{x}_t$ to the cell input, input gate, forget gate, and output gate, respectively. Similarly, the recurrent weights $r_{z}$, $r_{i}$, $r_{f}$, and $r_{o}$ facilitate the connections between the previous hidden state $h_{t-1}$ and the respective cell input, input gate, forget gate, and output gate. Bias terms associated with each of these are $b_{z}$, $b_{i}$, $b_{f}$, and $b_{o}$. The activation functions $\varphi$ and $\psi$ correspond to the cell input and hidden state, with $\tanh$ most commonly used for both. The role of $\psi$ is specifically to bound the cell state, which could otherwise grow without limit. All gates use the sigmoid function, $\sigma \left( x \right)= 1/(1+\exp(-x))$, as their activation. 

Subsequent developments consolidated individual scalar memory cells into a vector representation $\mathbf{c}_t \in \mathbb{R}^{d}$, which enables the introduction of recurrent weight matrices $\mathbf{R} \in \mathbb{R}^{d \times d}$ to permit the intermixing of outputs across multiple memory cells~\citep{Greff:15} (for comprehensive details, refer to Appendix~\ref{sec:apporgLSTM}). Empirical ablation studies have confirmed that each memory cell component plays an essential role in overall performance~\citep{Greff:15}.

\subsection{sLSTM}
\label{sec:sLSTM}

To enable LSTMs to adjust their storage choices, we incorporate exponential gates (shown in red), complemented by normalization and stabilization mechanisms. Specifically, the activation functions for both the input and forget gates may be exponential. For normalization purposes, we implement a normalizer state, which accumulates the products of the input gate and each subsequent forget gate.

The scalar sLSTM forward pass is:

\begin{align}
\label{eq:slstmforward}
c_t \ &= \  f_t \ c_{t-1} \ + \ i_t \ z_t & & 
 &\text{cell state} \\
n_t \ &= \  f_t \ n_{t-1} \ + \ i_t  & & 
 &\text{normalizer state} \\
h_t \ &= \ o_t \ \tilde{h}_t \ ,
  &\tilde{h}_t \ &= \ c_t / n_t
&\text{hidden state} \\
z_t \ &= \ \varphi \left( \tilde{z}_t \right) \ , 
  &\tilde{z}_t \ &=  \ \mathbf{w}^\top_{z} \ \mathbf{x}_t \ + \
  r_{z}  h_{t-1} \ + \  b_{z}
  &\text{cell input} \\
i_t \ &= \ \!\exp \left( \tilde{i}_t  \right) \ , 
  &\tilde{i}_t \ &= \ \mathbf{w}^\top_{i} \ \mathbf{x}_t \ + \
  r_{i}  \ h_{t-1} \ + \  b_{i}
  &\text{input gate} \\
f_t \ &= \ \sigma \left(  \tilde{f}_t \right) \ \text{OR} \
\!\exp \left(  \tilde{f}_t \right) \ ,
  &\tilde{f}_t \ &= \ \mathbf{w}^\top_{f} \ \mathbf{x}_t  \ + \
  r_{f}  \ h_{t-1} \ + \  b_{f}
  &\text{forget gate} \\
o_t \ &= \ \sigma \left( \tilde{o}_t \right) \ , 
  &\tilde{o}_t  \ &= \ \mathbf{w}^\top_{o} \ \mathbf{x}_t \ + \
  r_{o}  \ h_{t-1} \ + \  b_{o}
  &\text{output gate} 
\end{align}

We transfer the original LSTM gating techniques, i.e., input- and/or hidden-dependent gating plus bias term, to the new architectures. Exponential activation functions can lead to large values that cause overflows. Therefore, we stabilize gates with an additional state \mbox{$m_t$~\citep{Milakov:18arxiv}}:

\begin{align}
\label{eq:slstmstabil}
m_t \ &= \ \max \left( \log ( f_t ) + m_{t-1} , \log ( i_t ) \right) &\text{stabilizer state} \\
i'_t \ &= \ \!\exp \left( \log \left ( i_t \right) - m_t \right) \ = \ \!\exp \left( \tilde{i}_t - m_t \right) \ \, 
  &\text{stabil. input gate} \\
f'_t \ &= \ \!\exp \left( \log \left( f_t \right) + m_{t-1} - m_t \right) \ \, 
  &\text{stabil. forget gate}
\end{align}

As demonstrated in Appendix~\ref{sec:appsLSTM}, substituting $f_t$ with $f'_t$ and $i_t$ with $i'_t$ during the forward computation leaves both the network's overall output and the gradients of the loss with respect to the parameters unchanged.

\paragraph{New Memory Mixing.} Similar to the standard LSTM, the sLSTM architecture supports multiple memory cells (refer to Appendix~\ref{sec:appsLSTM}). The presence of several memory cells allows for the mixing of memory via the recurrent matrices $\mathbf{R}_{\mathbf{z}}$, $\mathbf{R}_{\mathbf{i}}$, $\mathbf{R}_{\mathbf{f}}$, $\mathbf{R}_{\mathbf{o}}$, which connect the hidden state vector $\mathbf{h}$ to the memory cell input $\mathbf{z}$, as well as to the gating signals $\mathbf{i}$, $\mathbf{f}$, and $\mathbf{o}$. What distinguishes this approach to memory mixing is the incorporation of exponential gating. This newer sLSTM variant can organize its architecture into multiple heads, facilitating intra-head memory mixing but excluding inter-head mixing. The combination of sLSTM heads with exponential gating introduces an innovative mechanism for memory integration.

\subsection{mLSTM}
\label{sec:mLSTM}

To improve the memory capacity of LSTMs, the memory cell is expanded from a single scalar $c \in \mathbb{R}$ to a matrix $\mathbf{C} \in \mathbb{R}^{d \times d}$. As a result, information retrieval now involves matrix multiplication. At each time step $t$, the objective is to memorize a pair of vectors: the key $\mathbf{k}_t \in \mathbb{R}^d$ and the value $\mathbf{v}_t \in \mathbb{R}^d$, in accordance with the terminology adopted from Transformer models. Subsequently, at time $t + \tau$, the value $\mathbf{v}_t$ is expected to be accessed using a query vector $\mathbf{q}_{t+\tau} \in \mathbb{R}^d$. This framework aligns with the principles of Bidirectional Associative Memories (BAMs) \citep{Kohonen:72,Anderson:72,Nakano:72,Anderson:77}.

The covariance update rule~\citep{Sejnowski:77,Dayan:91} for storing a key-value pair is

\begin{align}
\mathbf{C}_t \ &= \ \mathbf{C}_{t-1} \ + \ \mathbf{v}_{t} \ \mathbf{k}_t^\top \ .
\end{align}

We posit the application of layer normalization prior to the projection of inputs into keys and values, ensuring that these vectors possess zero mean. The rule for updating covariance is provably optimal~\citep{Dayan:91} when the objective is to maximize the distinguishability of retrieved binary vectors; this is tantamount to achieving the highest possible signal-to-noise ratio. If retrieval is restricted to pairwise interactions—and one accepts the quadratic computational cost as seen in attention mechanisms~\citep{Krotov:16,Krotov:17,Ramsauer:21}—an even greater degree of separability becomes achievable. Notably, the covariance update process aligns with the Fast Weight Programmers framework \citep{Schmidhuber:92ncfastweights,Schlag:21}, wherein modifications have introduced a fixed decay factor applied to $\mathbf{C}_{t-1}$ and a fixed learning rate applied to the 
$\mathbf{v}_{t} \mathbf{k}_t ^\top$ component~\citep{Ba:16}. Following this perspective, we embed the covariance update rule within the LSTM architecture: the forget gate serves as a decay parameter, the input gate governs the learning rate, and the output gate modulates the magnitude of the retrieved vector.

For this matrix memory, the normalizer state is the weighted sum of key vectors, where each key vector is weighted by the input gate and all future forget gates. Again, the normalizer state keeps record of the strength of the gates. Since the dot product between query and normalizer state can be close to zero, we use the absolute value of this dot product and lower bound it by a threshold (typically 1.0) as done previously~\citep{Sun:23arxiv}. The mLSTM forward pass is:

\begin{align}
\label{eq:mlstm_recurrent_begin}
\mathbf{C}_t \ &= \  f_t \ \mathbf{C}_{t-1} \ + \ 
  i_t \ \mathbf{v}_t \ \mathbf{k}_t^\top & &  &\text{cell state} \\
\mathbf{n}_t \ &= \  f_t \ \mathbf{n}_{t-1} \ + \ 
  i_t \ \mathbf{k}_t & &  &\text{normalizer state} \\
\mathbf{h}_t  \ &= \ \mathbf{o}_t \ \odot \ \tilde{\mathbf{h}}_t
\ , \qquad \qquad 
\tilde{\mathbf{h}}_t \ = \   \mathbf{C}_t \mathbf{q}_t \ / \ 
  \max \left\{ |\mathbf{n}_t^\top \mathbf{q}_t|, 1 \right\} 
   &&&\text{hidden state} \\
\mathbf{q}_t \ &= \ \mathbf{W}_q \ \mathbf{x}_t \ + \ \mathbf{b}_q  & &  &\text{query input} \\
\mathbf{k}_t \ &= \ \frac{1}{\sqrt{d}} \mathbf{W}_k \ \mathbf{x}_t \ + \ \mathbf{b}_k  & &  &\text{key input} \\
\mathbf{v}_t \ &= \ \mathbf{W}_v \ \mathbf{x}_t \ + \ \mathbf{b}_v  & &  &\text{value input} \\
i_t \ &= \ \!\exp \!  \! \left( \tilde{i}_t  \right) \ , 
 \qquad \qquad \ \ \ \ \,
  \tilde{i}_t \ = \ \mathbf{w}^\top_{i} \ \mathbf{x}_t \ + \  b_{i} \ \ \ 
  &&&\text{input gate} \\
f_t \ &= \ \sigma \!  \left(  \tilde{f}_t \right) \ \text{OR} \ \!\exp \!  \! \left(  \tilde{f}_t \right) , \ \ \: \!
\tilde{f}_t \ = \ \mathbf{w}^\top_{f} \ \mathbf{x}_t  \ + \
  b_{f}
  &&& \text{forget gate} \\
\label{eq:mlstm_recurrent_end}
\mathbf{o}_t \ &= \ \sigma \left( \tilde{\mathbf{o}}_t \right) \ , \qquad \qquad \quad \ \ \ \,
  \tilde{\mathbf{o}}_t  \ = \ \mathbf{W}_{\mathbf{o}} \ \mathbf{x}_t \ + \
  \mathbf{b}_{\mathbf{o}}
  &&&\text{output gate} 
\end{align}

mLSTM accommodates several memory cells similarly to standard LSTM architectures. Within mLSTM, the concepts of multiple heads and multiple cells function equivalently due to the absence of memory integration across cells. To ensure the stability of mLSTM’s exponential gating mechanisms, the stabilization procedures implemented for sLSTM are also applied here (refer to Equation~\ref{eq:slstmstabil}).  
Because memory mixing does not occur in mLSTM, the recurrence relation can be re-expressed in a parallelized format. Further specifics are available in Appendix~\ref{sec:appmLSTM}.

\subsection{xLSTM Architecture}
\label{sec:xLSTMarch}

\paragraph{xLSTM Blocks.}
\label{sec:xLSTMblock}
The function of an xLSTM block is to capture the sequence’s historical information within a high-dimensional, nonlinear representation. This approach aids in disentangling distinct past contexts, an essential condition for accurately inferring the forthcoming element in a sequence—such as predicting the next token. The motivation stems from Cover’s Theorem~\citep{Cover:65}, which asserts that when patterns are mapped non-linearly into a space of increased dimensionality, they are more readily separable using linear boundaries compared to the original, lower-dimensional space.

We examine two structural variants of the residual block: (i) the “post up-projection” design (analogous to that of Transformers), which produces a nonlinear summary in the native dimension before elevating the representation through a linear projection to a higher-dimensional space. Subsequently, a nonlinear activation is applied, followed by a linear projection that returns the representation to the original dimension. This is illustrated in the left section of Figure~\ref{fig:Backbones} and the third column of Figure~\ref{fig:xlstm_sketch}, with a more granular depiction available in Figure~\ref{fig:appsLSTM_detailed} of the appendix. (ii) The “pre up-projection” configuration (resembling mechanisms in State Space Models), whereby the input undergoes an initial linear transformation into a higher-dimensional space,

compresses historical information in a non-linear fashion within a high-dimensional latent space, followed by a linear transformation that restores information to the original space. In the architecture, an xLSTM block that utilizes an sLSTM incorporates the post up-projection module. Conversely, when an xLSTM block includes an mLSTM, the pre up-projection is employed, taking advantage of the enlarged memory capacity characteristic of high-dimensional representations. For further illustration, please consult the left side of Figure~\ref{fig:Backbones}, the third column in Figure~\ref{fig:xlstm_sketch}, or Figure~\ref{fig:appsLSTM_detailed} provided in the appendix.

\paragraph{xLSTM Architecture. }
\label{sec:xLSTMarchitecure}
The xLSTM model is built through the residual composition of modular blocks~\citep{Srivastava:15,He:16}. Its foundation is the widely adopted pre-LayerNorm~\citep{Ba:16b} residual backbone, which has become the standard in recent Large Language Model designs. Further details are depicted in the final column of Figure~\ref{fig:xlstm_sketch}.

\subsection{Memory and Speed Considerations}

In contrast to Transformers, xLSTM architectures exhibit both linear computational cost and fixed memory requirements relative to sequence length. Because xLSTM employs compressive memory, it is particularly advantageous for use in industrial scenarios and for deployment at the edge.

While mLSTM’s memory mechanism does not rely on additional parameters, it does introduce significant computational expense due to the involvement of a $d \times d$ matrix both for memory storage and its updates. This creates a trade-off between increased memory capacity and greater computational demands. Despite this, the associated calculations are amenable to parallelization on GPUs, which mitigates their practical impact on total processing time.

Although mLSTM allows for parallel computation similarly to approaches such as FlashAttention~\citep{Dao:22,Dao:23} or GLA~\citep{Yang:23arxiv}, sLSTM lacks this capability because of its memory mixing introduced by hidden-to-hidden connections. Nonetheless, we have introduced an optimized CUDA-based implementation for sLSTM, utilizing GPU memory all the way to the register level. As a result, our optimized sLSTM is usually less than twice as slow as mLSTM in practice.

\section{Related Work}
\label{sec:related}

\paragraph{Linear Attention.} To address the quadratic scaling of the Transformer's attention with respect to context length, various strategies have been proposed to achieve linear complexity. The Synthesizer approach generates learned synthetic attention weights independently of token-to-token dependencies~\citep{Tay:20arxiv}. Linformer employs a low-rank projection to approximate self-attention in a manner that scales linearly~\citep{Wang:20arxiv}. Linear Transformer redefines the attention operation so that its computational requirements are linear~\citep{Katharopoulos:20}. Performer leverages positive orthogonal random features to provide a linear estimation of the softmax attention mechanism~\citep{Choromanski:21}. Additionally, Structured Global Convolution (SGConv)~\citep{Li:22} and the Hyena Hierarchy~\citep{Poli:23} have replaced traditional attention with rapid long-range convolutional operations.

\paragraph{State Space Models.} State Space Models (SSMs) have gained significant traction in recent work, chiefly due to their linear scaling with sequence length and competitive performance relative to Transformer architectures. Early developments in this area include the Structured State Space sequence model (S4)~\citep{Gu:21}. Subsequently, various extensions have been proposed, such as the Diagonal State Space (DSS) model~\citep{Gupta:22}, Gated State Space (GSS) models~\citep{Mehta:22}, S5 model~\citep{Smith:22}, Bidirectional Gated SSM (BiGS)~\citep{Wang:22}, H3 model~\citep{Fu:23}, and Mamba~\citep{Gu:24arxiv}.

\paragraph{Recurrent Neural Networks.} In recent literature, Recurrent Neural Networks (RNNs) have been advocated as alternatives to Transformer architectures and attention mechanisms, primarily owing to their computational linearity with respect to sequence length. Deep Linear Recurrent Units (LRUs) within RNNs have demonstrated favorable outcomes in language modeling tasks~\citep{Orvieto:23, De:24arxiv}. Similarly, Hierarchically Gated Linear RNN (HGRN)~\citep{Qin:23} and its successor HGRN2~\citep{Qin:24arxiv} have also produced encouraging results. Among RNN approaches tailored for large-scale language modeling, RWKV~\citep{Peng:23arxivshort,Peng:24arxivshort} stands out, achieving performance levels comparable to Transformer models.

\paragraph{Gating.}
LSTM's incorporation of gating mechanisms has served as a foundational concept, inspiring many modern methods that reinterpret or extend the gating idea. Such gating structures have been incorporated into various architectures including HGRN~\citep{Qin:23}, HGRN2~\citep{Qin:24arxiv}, Gated Linear Attention (GLA)~\citep{Yang:23arxiv}, Gated State Space (GSS) models~\citep{Mehta:22}, Bidirectional Gated SSM (BiGS)~\citep{Wang:22}, Moving Average Equipped Gated Attention (MEGA)~\citep{Ma:22}, RWKV~\citep{Peng:23arxivshort}, as well as Mamba~\citep{Gu:24arxiv}.

\paragraph{Covariance Update Rule.}
In an effort to boost memory capacity, the mLSTM cell was augmented with a matrix memory utilizing a covariance update strategy. This approach aligns with techniques employed in other models such as Fast Weight Programmers \citep{Schmidhuber:92ncfastweights,Schlag:21}, RWKV-5 and RWKV-6~\citep{Peng:24arxivshort}, Retention~\citep{Sun:23arxiv}, Linear Transformer~\citep{Katharopoulos:20}, and HGRN2~\citep{Qin:24arxiv}.

\paragraph{Most Related.} The models most akin to xLSTM from a conceptual standpoint include Retention~\citep{Sun:23arxiv}, RWKV~\citep{Peng:23arxivshort,Peng:24arxivshort}, and HGRN2~\citep{Qin:24arxiv}. These approaches all utilize the principles of matrix memory and/or gating mechanisms. Nonetheless, unlike the newly proposed sLSTM, these methods do not support memory mixing. The presence of memory mixing in LSTMs facilitates effective state tracking, a crucial capability for tasks such as code evaluation or monitoring entities throughout extended narratives. As a result, LSTMs demonstrate greater expressive power than both State Space Models (SSMs) and Transformers~\citep{Merrill:24,Deletang:23} due to their ability to handle state tracking challenges.

\paragraph{Residually Stacking Architectures.} The xLSTM architectures, consistent with the prevailing approach for recent large-scale deep learning models, are assembled by layering building blocks using residual connections~\citep{Srivastava:15,He:16}.

This design has proven instrumental in facilitating the development of deep convolutional neural networks~\citep{He:16} and Transformer-based models~\citep{Vaswani:17}. Transformers have become the foundational technology driving Large Language Models (LLMs) such as GPT-3~\citep{Brown:20short}, ChatGPT~\citep{Schulman:22short}, GPT-4~\citep{Achiam:23gpt4short}, Megatron-LM~\citep{Shoeybi:19}, Gopher~\citep{Rae:21short}, ERNIE 3.0 Titan~\citep{Wang:21arxivshort}, GLaM~\citep{Du:21glamshort}, Chinese M6~\citep{Lin:21}, multilingual AlexaTM 20B~\citep{Soltan:22}, OPT~\citep{Zhang:22opt}, Chinchilla~\citep{Hoffmann:22short}, BLOOM~\citep{Scao:22arxivshort}, GLM-130B~\citep{Zeng:22arxivshort}, LaMDA~\citep{Thoppilan:22short}, PaLM~\citep{Chowdhery:22short}, Llama~\citep{Touvron:23arxiv}, and Gemini~\citep{GeminiTeam:23arxiv,Reid:24arxiv}.

\section{Experiments}
\label{sec:Exp}

This section presents the empirical analysis of xLSTM, placing particular emphasis on its performance in language modeling tasks. Section~\ref{sec:ExpSyntethic} explores xLSTM's behavior on artificial benchmarks, providing insights into its unique properties. Section~\ref{sec:ExpComparison} details a comparison of validation perplexity across a variety of state-of-the-art language models, all trained using 15B tokens drawn from SlimPajama~\citep{Soboleva:23}. In addition, we carry out ablation studies targeting xLSTM's architecture, and subsequently examine the scaling dynamics of each model, employing methodologies from \citet{Kaplan:20} and \citet{Brown:20short}. 

In Section~\ref{sec:ExpLanguage}, we broaden our language modeling evaluation. Here, we analyze xLSTM alongside the strongest models identified in Section~\ref{sec:ExpComparison} after training them with 300B SlimPajama tokens~\citep{Soboleva:23}. The analysis covers four aspects: (1) performance on length generalization to extended contexts, (2) validation perplexity as well as downstream task outcomes~\citep{Sutawika:24}, (3) results on the 571 domain subsets from the PALOMA benchmark dataset~\citep{Magnusson:23arxivshort}, and (4) the scaling properties of all methods, considering the expanded dataset that is 20 times larger than previously used.

\label{sec:xlstm_blocks_notation}
Throughout all experiments, we denote xLSTM[$a$:$b$] to indicate the proportion $a/b$ between xLSTM blocks built with mLSTM and those built with sLSTM. For instance, the notation xLSTM[7:1] specifies a configuration where seven out of eight blocks utilize mLSTM, while one block uses sLSTM. In a standard case of 48 total blocks, this setup would include 42 mLSTM blocks and 6 sLSTM blocks. Additionally, in all experimental setups, both mLSTM and sLSTM are paired with pre up-projection and post up-projection blocks, respectively.

\subsection{Synthetic Tasks and Long Range Arena}
\label{sec:ExpSyntethic}
Initially, we evaluate the capabilities of xLSTM’s novel exponential gating mechanism with memory mixing using formal languages~\citep{Deletang:23}. Subsequently, the new matrix memory feature in xLSTM is examined on the Multi-Query Associative Recall task~\citep{Arora:23arxiv}. Lastly, xLSTM’s handling of extended sequence inputs is analyzed through its performance on the Long Range Arena benchmark~\citep{Tay:21}.

\paragraph{Evaluating xLSTM's Exponential Gating with Memory Mixing.} To assess the capabilities of xLSTM's recently introduced exponential gating mechanism with memory mixing—designed to address state tracking tasks~\citep{Merrill:24, Merrill:23}—we extend the formal language benchmarks from \citet{Deletang:23} to incorporate multi-length sequences, which facilitates testing the model's ability to extrapolate to longer inputs. Details about the individual tasks and a comprehensive presentation of results are provided in Appendix~\ref{sec:appExpSynthetic-formalLang}. 

Our evaluation contrasts xLSTM against a range of approaches including Transformers, State Space Models, and various Recurrent Neural Networks. Accuracy is computed only on task-specific target tokens, with values normalized from 0 (indicating random chance) to 1 (indicating perfect predictions). We benchmark 2-block versions of the following architectures: xLSTM[0:1] (sLSTM-only), xLSTM[1:0] (mLSTM-only), xLSTM[1:1], Llama, Mamba, RWKV, Retention, Hyena, LSTM, and LSTM implemented within Transformer blocks (LSTM (Block)). Results are summarized in Figure~\ref{fig:formal-main}.

Architectures such as standard Transformers and State Space Models, when lacking the memory mixing mechanism necessary for state retention, are incapable of solving even relatively simple regular grammar tasks like parity. This observation is consistent with recent evidence indicating that, with respect to state tracking, Transformers and State Space Models are intrinsically weaker than RNN-based approaches~\citep{Merrill:24, Merrill:23, Deletang:23}.

\paragraph{Test of xLSTM's Memory Capacities on Associative Recall Tasks.} In this study, we investigate the effectiveness of xLSTM's newly introduced matrix memory by assessing its memory capacity using the Multi-Query Associative Recall task~\citep{Arora:23arxiv}. Sequences are constructed by randomly sampling key-value pairs from an extensive vocabulary, requiring the models to memorize them for subsequent retrieval. To increase the challenge beyond the original setup, we raise the number of key-value pairs to as many as 256, and lengthen the context window to 2048 tokens. This expansion enables a more comprehensive evaluation of each model's memory capacity. We examine the performance of 2-block versions of Llama, Mamba, RWKV-5, RWKV-6, xLSTM[1:1], and xLSTM[1:0], measuring their accuracy in correctly recalling stored key-value pairs. Transformers, such as Llama, are known for having memory that scales exponentially with the coding dimension~\citep{Ramsauer:21}, thereby serving as a performance benchmark for this task. Figure~\ref{fig:mqar-main} summarizes our main findings.
Among non-Transformer approaches, xLSTM[1:1] outperforms its peers, including on tasks with smaller model sizes. Notably, the integration of the sLSTM block does not impair the memory capabilities; instead, it appears to further enhance them, particularly evident in scenarios demanding the recall of 256 key-value pairs. Appendix~\ref{sec:appExpSynthetic-mqar} contains supplemental analyses, including extrapolation experiments suggesting that the improved memory mechanisms in xLSTM enable successful generalization to longer contexts than were encountered during training.

\paragraph{Test of xLSTM's Long Context Capabilities on Long Range Arena.} To evaluate the effectiveness of xLSTM in processing extended sequences and sizable contextual information, we benchmark various approaches using the Long Range Arena~\citep{Tay:21}. Across every task, xLSTM reliably achieves high performance, indicating that its architecture excels at managing the diverse challenges associated with long-context scenarios.

For more details, see Appendix~\ref{sec:appExpSynthetic-lra}.

\subsection{Method Comparison and Ablation Study}
\label{sec:ExpComparison}

In order to investigate the central inquiry of this work—namely, the performance of our novel LSTM variants when scaled for language modeling tasks—we conduct experiments in an auto-regressive framework. Specifically, we train xLSTMs, Transformers, State Space Models, and additional approaches using 15B tokens from SlimPajama. Subsequent evaluation is carried out on a validation dataset, alongside ablation analyses focusing on the xLSTMs.

\paragraph{Evaluating xLSTM Against Alternative Approaches.} To facilitate a fair assessment, each model is trained on 15B tokens sourced from SlimPajama~\citep{Soboleva:23}, with validation set perplexity serving as the evaluation metric. We benchmark xLSTM (our contribution) alongside multiple established methods: GPT-3 (Transformer)~\citep{Brown:20short}, Llama (Transformer)~\citep{Touvron:23arxiv}, H3 (SSM)~\citep{Fu:23}, Mamba (SSM)~\citep{Gu:24arxiv}, RWKV-4 (RNN)~\citep{Peng:23arxivshort}, RWKV-5 (RNN)~\citep{Peng:24arxivshort}, RWKV-6 (RNN)~\citep{Peng:24arxivshort}, GLA (linear Transformer)~\citep{Yang:23arxiv}, HGRN (RNN)~\citep{Qin:23}, HGRN2 (RNN)~\citep{Qin:24arxiv}, RetNet (linear Transformer)~\citep{Sun:23arxiv}, Hyena (linear Transformer)~\citep{Poli:23}, xLSTM[1:0], and xLSTM[7:1], as described in Section~\ref{sec:xlstm_blocks_notation}. Model training utilized mixed precision; however, for RWKV-5, RWKV-6, GLA, and HGRN2, mixed-precision execution was implemented without the automated PyTorch mixed precision tools (details in Appendix Section~\ref{sec:appGenTrainProc}). Our classification framework groups the models as (a) Transformers, (b) State Space Models (SSMs), and (c) Recurrent Neural Networks (RNNs) and linear Transformers, where linear Transformers replace the conventional Transformer attention with linear alternatives. All models are architecturally equivalent to a GPT-3 with 350M parameters, adopting an embedding size of 1024 and incorporating 24 residual blocks. GPT-3 uniquely employs shared token and output embedding weights, resulting in a reduced parameter count. As presented in Table~\ref{tab:spaj15b_model_results}, xLSTM achieves superior validation perplexity in comparison to all other tested methods. Further discussion is available in Appendix~\ref{sec:appExpComparison}. Scaling trends observed in Figure~\ref{fig:temp_sclaw15B} suggest that xLSTM maintains its advantageous performance as model size increases.

\paragraph{Ablation Studies.}
Table~\ref{tab:spaj15b_model_results} and Figure~\ref{fig:temp_sclaw15B} reveal that xLSTM delivers strong language modeling performance when trained on 15B tokens drawn from SlimPajama. To systematically dissect the improvements contributed by xLSTM over the standard LSTM, we progressively modify a traditional LSTM architecture to transition towards xLSTM. The first modification involves embedding LSTM layers within pre-LayerNorm residual structures. Next, we incorporate a post up-projection module. The final step is the addition of exponential gating and matrix memory mechanisms.

The results are shown in Table~\ref{tab:ablstudies} (top). 

Ablation experiments link the notable gains in performance to the combination of exponential gating and matrix memory. Furthermore, recognizing the critical role of gating within RNNs and State Space Models, we examine alternative gating strategies. As shown in Table~\ref{tab:ablstudies} (bottom), our findings indicate that assigning learnable parameters to each gate and making them responsive to input data further enhances performance incrementally. More comprehensive analyses concerning the backbone elements can be found in Appendix~\ref{sec:appAblDetails}.

\subsection{xLSTM as Large Language Model}
\label{sec:ExpLanguage}

In concluding this work, we conduct extensive language modeling experiments to evaluate xLSTM's viability as a large language model. We significantly expand the training dataset, utilizing 300B tokens sourced from SlimPajama, which aligns with the token count used in prior studies such as Mamba~\citep{Gu:24arxiv} and Griffin~\citep{De:24arxiv}.

For a comprehensive comparison, xLSTM is measured against RWKV-4, Llama, and Mamba—each representing the distinct methodological categories described in Section~\ref{sec:ExpComparison}. RWKV-4 is chosen to represent RNNs, as suitable training precision configurations for RWKV-5, RWKV-6, and HGRN2 were identified only after commencement of the 300B token training runs (refer to Appendix~\ref{sec:appGenTrainProc}).

Our experiments encompass various model scales (125M, 350M, 760M, and 1.3B parameters). We examine all models for their capacity to extrapolate sequence lengths, analyze their validation set performance, gauge effectiveness across downstream tasks, evaluate their language modeling ability across 571 distinct text domains using the PALOMA benchmark, and thoroughly study their scaling laws.

\paragraph{Sequence Length Extrapolation.}
\label{sec:spaj300B_ctx_extrapolation}
To begin, we assess how well the 1.3B-parameter versions of xLSTM, RWKV-4, Llama, and Mamba generalize to longer sequence lengths. All models are initially trained using sequences of 2048 tokens and subsequently evaluated on sequences extended up to 16384 tokens. The outcomes, depicted in Figure~\ref{fig:spaj300B_seq_extrapolation}, reveal that xLSTM continues to achieve low perplexity at increased context lengths, distinguishing it from alternative approaches.

\paragraph{Validation Perplexity and Downstream Tasks.}
\label{sec:spaj_downstream_eval}
Next, across every parameter scale, we benchmark xLSTM, RWKV-4, Llama, and Mamba on next token prediction over the SlimPajama validation split, as well as on downstream tasks intended to evaluate abilities in common sense reasoning.

In Table~\ref{tab:lmeval}, the third column reports the validation set perplexities associated with various approaches. Among these, xLSTM[1:0] and xLSTM[7:1] consistently achieve the lowest perplexity scores across every model size. Performance metrics for downstream tasks are presented in the remaining columns of Table~\ref{tab:lmeval}. Across nearly all evaluated tasks and model sizes, xLSTM demonstrates superior performance, with the exception of the ARC task where Mamba outperforms xLSTM in select cases. Additional information can be found in Appendix~\ref{sec:appExpLanguage}.

\paragraph{Performance on PALOMA Language Tasks.} As a third evaluation, we assess the next token prediction accuracy across all model scales for xLSTM, RWKV-4, Llama, and Mamba using PALOMA language tasks~\citep{Magnusson:23arxivshort}. Perplexity is utilized as the metric to quantify next token prediction across 571 diverse text domains, which encompass sources like nytimes.com and Reddit's r/depression.

Table~\ref{tab:ppleval} summarizes the perplexity results, partitioned into language modeling (first seven columns) and specialized domain benchmarks (final five columns). Notably, xLSTM[1:0] surpasses xLSTM[7:1] in language task performance.

xLSTM[1:0] exhibits lower perplexity compared to Mamba in 568 of 571 (99.5\%) domains, outperforms Llama in 486 of 571 (85.1\%) domains, and achieves better results than RWKV-4 in 570 of 571 (99.8\%) domains. Additional data can be found in Appendix~\ref{sec:appExpLanguage}.

\paragraph{Scaling Laws.} Next, we investigate the scaling properties characterized by power-law relationships, which facilitate the projection of model performance as size increases~\citep{Kaplan:20,Brown:20short}. Figure~\ref{fig:temp_sclaw300B} depicts this scaling trend. Although all evaluated models exhibit broadly analogous scaling patterns, the baselines are offset from one another. Among them, RWKV-4 yields the lowest performance, with Llama and Mamba trailing behind. Notably, xLSTM outperforms Mamba, mirroring the performance gap observed between Mamba and Llama. This scaling trend suggests that xLSTM's relative advantage over Transformers and State-Space models persists as model capacity grows.

\textbf{Generation Times and Maximal Throughput.} We report measurements of text generation latency in Figure~\ref{fig:inference_time_generation} and the peak throughput in Figure~\ref{fig:inference_time_decoding_throughput} (left) for our various xLSTM configurations, all at the 1.3B parameter scale. For context, we compare these results with comparable Mamba, Llama, and RWKV implementations provided by HuggingFace, incorporating a static key-value cache in the Llama setup. Notably, at the time these experiments were conducted, compiling the full cache for the Transformer Model and compiling the Mamba implementation using \texttt{torch.compile} was not feasible. In all text generation tests, models are evaluated at a batch size of 1 and a pre-fill of 16, a setting particularly advantageous for the Transformer architecture.

As illustrated in Figure~\ref{fig:inference_time_generation}, both xLSTM and fellow recurrent architectures such as Mamba and RWKV-4 demonstrate linear computational scaling, whereas Llama exhibits quadratic scaling. For the decoding throughput experiments, varying batch sizes and pre-fill values are explored for Llama. The data in Figure~\ref{fig:inference_time_decoding_throughput} (right) demonstrates that, by virtue of its constant memory usage, xLSTM accommodates significantly higher batch sizes relative to Llama, thereby obtaining superior throughput.

\section{Limitations}

(i) Unlike mLSTM, the memory mixing in sLSTM restricts operations that can be efficiently parallelized, making rapid parallel implementations unfeasible. However, we engineered a specialized CUDA kernel for sLSTM, which currently operates at less than twice the latency of our parallel mLSTM version. (ii) The mLSTM CUDA kernels lack optimization at present, resulting in a performance that is approximately four times slower than FlashAttention or the scanning approach utilized in Mamba. Enhanced CUDA kernels could be developed following the methodology of FlashAttention. (iii) Managing the matrix memory in mLSTM incurs substantial computational demands because it involves $d \times d$ matrices. Nonetheless, since the memory update and access operations are non-parametric and amenable to parallel processing via common matrix computations, the actual increase in wall clock time due to complex memory remains limited. (iv) Careful selection of initial values for the forget gates is essential. (v) As the matrix memory does not scale with sequence length, extending the sequence might challenge memory capacity for lengthy contexts; however, this does not present a constraint for sequence contexts up to 16k, as discussed in Section~\ref{sec:spaj300B_ctx_extrapolation}. (vi) Owing to the significant computational requirements of large-scale language modeling, we have not thoroughly optimized either the model architecture or the hyperparameters, particularly for larger xLSTM models. We recognize that rigorous optimization will be necessary for xLSTM architectures to realize their full capabilities.

\section{Conclusion}
We have addressed, to an extent, the question: How effective does language modeling become when LSTM architectures are scaled to the billions of parameters? Up to this point, our findings indicate: ``At minimum, their performance rivals that of prevailing techniques such as Transformers and State Space Models''. We introduced xLSTM, which augments the original LSTM via exponential gating combined with memory mixing and a revised memory architecture. On language modeling tasks, xLSTM models demonstrate competitive or superior results when assessed against leading methods, notably Transformers and State Space Models. Observed scaling laws suggest that expanding xLSTM models further may position them as strong contenders to current Transformer-based Large Language Models. Moreover, xLSTM could play an influential role in domains beyond language modeling, including Reinforcement Learning, Time Series Prediction, and the simulation of physical systems.

\end{document}